\documentclass{article}
%  ╔══════════════════════════════════════════════════════════╗
%  ║                         Title                            ║
%  ╚══════════════════════════════════════════════════════════╝
%NOTE: I want \title,\author to be near top, and this \renewcommand{\maketitle} fails if these are above it.
\usepackage[utf8]{inputenc}         % Used to compile any UTF-8 Character (and supports Unicode)
\usepackage{titling}

\renewcommand{\maketitle}{
    \quad
    \vspace{1em}
  \begin{center}
      \LARGE\bfseries \thetitle \par
      \vspace{0.6em}
      \large
  \end{center}
    \vspace{1em}
}

\title{A Foray into the Fourier-Mukai Transform}
\author{Nathanael Chwojko-Srawley}
\date{\today}

%  ╔══════════════════════════════════════════════════════════╗
%  ║                         packages                         ║
%  ╚══════════════════════════════════════════════════════════╝

\usepackage{extarrows}         % Used to compile any UTF-8 Character (and supports Unicode)
\usepackage{stackengine}
\usepackage{colortbl}
\usepackage{scalerel}
\usepackage{amsmath, amssymb} % Used to write better mathematical expressions (ex. \mathbb{R})
\usepackage{stackrel}         % for left-right arrows, staking symbol above and bellow
\usepackage{mathrsfs}
\usepackage{bbm}                     % for mathbbm{1}
\usepackage{euscript}
\usepackage{commutative-diagrams}
\usepackage[font=itshape]{quoting}
\usepackage{mathtools}
  \usepackage{enumitem}
\usepackage{amsthm}           % Used to create theorem enviroments.
\usepackage{graphicx}         % Let's you use images in LaTeX; ex the \includegraphics command.
\usepackage{hhline}
%\usepackage{luatexja}
%\usepackage{fontspec}
\usepackage{dsfont}
\usepackage{wasysym}          % emoji's!
\usepackage{float}            % package with ``H'' for figures and tables, ex. Images, tables, etc
%\usepackage{subcaption}      % More options with sub-captions.
\usepackage{setspace}         % More flexibility on spacing in document.
\usepackage{hyperref}         % let's me put hyperlinks
%\usepackage{tikz}             % for commutative diagrams https://tex.stackexchange.com/questions/419221/how-to-do-the-pushout-with-universal-property
\usepackage{tikz-cd}          % for commutative diagrams https://tex.stackexchange.com/questions/419221/how-to-do-the-pushout-with-universal-property
%\usepackage{quiver}           % for better commutative diagrams
\usepackage{bookmark}         % creating heading shortcut bar on the left as in word
\usepackage{longtable}        % create tables that extend past one page
%\usepackage{siunitx}         % required for alignment
%\usepackage{multirow}        % for multiple rows
\usepackage{microtype}        % makes nice text. Fixes some under-/over- filled box problems.
%\usepackage{fancyvrb}        % Let's you write code (different style than listings).
\usepackage[most]{tcolorbox} 	      % Makes nice boxes for thms, cor, prop, or anything else.
%\usepackage{enumitem}        % More flexibility on enumeration (don't know much).
\usepackage{xcolor}           % Colour text.
\usepackage{wrapfig}          % To wrap text around figure
\usepackage{listings}         % Create nice colour code blocks (customized in preamble)
\usepackage{thmtools}         % Customize theorem environment (in preamble).
%\usepackage{marvosym}        % Provide ``basic'' symbols (hazard, religious, smiley face etc.)
\usepackage{array}            % \begin{table}{>{\em}c c} -> 1st column italic (bfseries for bold)
%\usepackage[english]{babel}  % If I want to blindtext in english.
\usepackage{blindtext}
\usepackage[a4paper, total={6in, 8.5in}]{geometry}
%\usepackage[symbol]{footmisc} % For daggers and asterisk for footnotes. 1: *, 2: dagger, etc.
\usepackage{fancyhdr}
\usepackage{titlesec}
%\usepackage{pst-plot}          %To make nice plots: https://tex.stackexchange.com/questions/47594/how-can-i-make-a-graph-of-a-function#47596
% \usepackage[answerdelayed]{exercise}
\usepackage{etoolbox}          % for Dynkin diagrams
\usepackage{dynkin-diagrams}   % for Dynkin diagrams

% ╔═════════════════════════════════════════════════════╗
% ║           for figures via inkspace                  ║
% ╚═════════════════════════════════════════════════════╝

% to import figures via: https://github.com/gillescastel/inkscape-figures
\usepackage{import}
\usepackage{pdfpages}
\usepackage{transparent}


% This command is the ONLY command imported via the packages snippet, think of it as a tiny package written out
\newcommand{\incfig}[2][1]{%
    \def\svgwidth{#1\columnwidth}
    \import{./figures/}{#2.pdf_tex}
}

% ╔═══════════════════════════════════════════════════════════╗
% ║                 bibliography and citing                   ║
% ╚═══════════════════════════════════════════════════════════╝

\usepackage{csquotes}
% \usepackage[backend=biber,citestyle=alphabetic,bibstyle=authortitle]{biblatex}
\usepackage{biblatex}

\addbibresource{bibliography.bib}

% ╔═══════════════════════════════════════════════════════════╗
% ║                    colours (colors)                       ║
% ╚═══════════════════════════════════════════════════════════╝

\definecolor{dkgreen}{rgb}{0,0.6,0}
\definecolor{gray}{rgb}{0.5,0.5,0.5}
\definecolor{mauve}{rgb}{0.58,0,0.82}
\definecolor{lightyellow}{rgb}{1,1,0.6}
\definecolor{reddish}{HTML}{A01A3A}

%  ╔══════════════════════════════════════════════════════════╗
%  ║                       book format                        ║
%  ╚══════════════════════════════════════════════════════════╝

\pagestyle{fancy}

\setlength\headheight{20pt}

\renewcommand{\sectionmark}[1]{\markright{\thesection.\ #1}}

\fancyfoot[C,CO]{\textcolor{gray}{Nathanael Chwojko-Srawley}}
\fancyfoot[R,RO]{\thepage}{}

% Create a new page style for the first page
\fancypagestyle{firstpage}{
  \fancyhf{} % Clear header and footer
  \fancyhead[L]{\theauthor}
  \fancyhead[R]{\today} % Date in the top-right
  \fancyfoot[C]{\textcolor{gray}{Nathanael Chwojko-Srawley}}
  \fancyfoot[R]{\thepage}
}

% Apply the first page style conditionally
\AtBeginDocument{\thispagestyle{firstpage}}

%have format the title command
\newcommand{\chapter}[1]{%
  \clearpage
  \vspace*{30pt} % Space before the chapter title
  {\normalfont\huge\bfseries #1\par} % Chapter title formatting
  \vspace{20pt} % Space after the chapter title
}


\setlength{\parindent}{0pt} % don't like indentation infront of paragraphs
\setlength{\parskip}{1ex}   %so that there is still space between paragarphs


%  ╔══════════════════════════════════════════════════════════╗
%  ║                       Shortcuts                          ║
%  ╚══════════════════════════════════════════════════════════╝

%\ExerciseLevelInToc

% I like original mathcal better, but need lowercase.
\DeclareFontFamily{OT1}{pzc}{}
\DeclareFontShape{OT1}{pzc}{m}{it}{<-> s * [1.10] pzcmi7t}{}
\DeclareMathAlphabet{\mathpzc}{OT1}{pzc}{m}{it}


% mathbb
\newcommand{\A}{{\mathbb{A}}}
\newcommand{\C}{{\mathbb{C}}}
\newcommand{\F}{{\mathbb{F}}}
\newcommand{\T}{{\mathbb{T}}}
\newcommand{\N}{{\mathbb{N}}}
\newcommand{\Q}{{\mathbb{Q}}}
\newcommand{\R}{{\mathbb{R}}}
\newcommand{\V}{{\mathbb{V}}}
\newcommand{\Z}{{\mathbb{Z}}}
\newcommand{\BI}{{\mathbb{I}}}
\renewcommand{\H}{{\mathbb{H}}}
\renewcommand{\P}{{\mathbb{P}}}


% mathcal
% \newcommand{\co}{{\mathfrak{o}}} %mathcal{o} looks like wreath product, so I changed it to mathfrak
\newcommand{\co}{{\mathpzc{o}}}
\newcommand{\CA}{{\mathcal{A}}}
\newcommand{\CB}{{\mathcal{B}}}
\newcommand{\CC}{{\mathcal{C}}}
\newcommand{\CD}{{\mathcal{D}}}
\newcommand{\CF}{{\mathcal{F}}}
\newcommand{\CG}{{\mathcal{G}}}
\newcommand{\CH}{{\mathcal{H}}}
\newcommand{\CI}{{\mathcal{I}}}
\newcommand{\CK}{{\mathcal{K}}}
\newcommand{\CL}{{\mathcal{L}}}
\newcommand{\CM}{{\mathcal{M}}}
\newcommand{\CN}{{\mathcal{N}}}
\newcommand{\CO}{{\mathcal{O}}}
\newcommand{\CE}{{\mathcal{O}}}
\newcommand{\CQ}{{\mathcal{Q}}}
\newcommand{\CR}{{\mathcal{R}}}
\newcommand{\CS}{{\mathcal{S}}}
\newcommand{\CT}{{\mathcal{T}}}
\newcommand{\CV}{{\mathcal{V}}}
\newcommand{\CZ}{{\mathcal{Z}}}
% EXCPETION:
\newcommand{\CP}{{\mathcal{P}}}
% \newcommand{\CP}{\mathbb{C}\mathrm{P}}
\newcommand{\RP}{\mathbb{R}\mathrm{P}}


%Euscript
\newcommand{\EB}{\EuScript{B}}
\newcommand{\EC}{{\EuScript{C}}}
\newcommand{\EF}{{\EuScript{F}}}
\newcommand{\EL}{{\EuScript{L}}}
\newcommand{\EO}{{\EuScript{O}}}


%mathfrak (lowercase)
\newcommand{\fa}{{\mathfrak{a}}}
\newcommand{\fb}{{\mathfrak{b}}}
\newcommand{\fc}{{\mathfrak{c}}}
\newcommand{\fd}{{\mathfrak{d}}}
\newcommand{\fm}{{\mathfrak{m}}}
\newcommand{\fn}{{\mathfrak{n}}}
\newcommand{\fo}{{\mathfrak{o}}}
\newcommand{\fp}{{\mathfrak{p}}}
\newcommand{\fq}{{\mathfrak{q}}}


%mathfrak (upper case)
\newcommand{\FA}{{\mathfrak{A}}}
\newcommand{\FB}{{\mathfrak{B}}}
\newcommand{\FC}{{\mathfrak{C}}}
\newcommand{\FD}{{\mathfrak{D}}}
\newcommand{\FK}{{\mathfrak{K}}}
\newcommand{\FM}{{\mathfrak{M}}}
\newcommand{\FN}{{\mathfrak{N}}}
\newcommand{\FO}{{\mathfrak{O}}}
\newcommand{\FP}{{\mathfrak{P}}}


%mathscr
\newcommand{\ff}{{\mathscr{F}}}
\newcommand{\SA}{\mathscr{A}}
\newcommand{\SC}{\mathscr{C}}
\newcommand{\SF}{\mathscr{F}}
\newcommand{\SG}{\mathscr{G}}
\newcommand{\SH}{\mathscr{H}}
\newcommand{\SI}{\mathscr{I}}
\newcommand{\SK}{\mathscr{K}}
\newcommand{\SN}{\mathscr{N}}
\newcommand{\SQ}{\mathscr{Q}}
\newcommand{\SR}{\mathscr{R}}
\newcommand{\SV}{\mathscr{V}}
\newcommand{\SZ}{\mathscr{Z}}


% operatorname -- 2 letters (lowercase and upper case)
\newcommand{\ab}{{\operatorname{ab}}}
\newcommand{\ad}{{\operatorname{ad}}}
\newcommand{\ch}{{\operatorname{ch}}}
\newcommand{\ev}{{\operatorname{ev}}}
\newcommand{\id}{{\operatorname{id}}}
\newcommand{\im}{{\operatorname{im}}}
\newcommand{\nm}{{\operatorname{Nm}}}
\newcommand{\op}{{\operatorname{op}}}
\newcommand{\rk}{{\operatorname{rk}}}
\newcommand{\sh}{{\operatorname{sh}}}
\newcommand{\tr}{{\operatorname{tr}}}

\newcommand{\Ab}{{\operatorname{Ab}}}
\newcommand{\Br}{{\operatorname{Br}}}
\newcommand{\Ch}{{\operatorname{Ch}}}
\newcommand{\Cl}{{\operatorname{Cl}}}
\newcommand{\GL}{{\operatorname{GL}}}
\newcommand{\Nm}{{\operatorname{Nm}}}
\newcommand{\SL}{{\operatorname{SL}}}

\renewcommand{\Re}{{\operatorname{Re}}}
\renewcommand{\Im}{{\operatorname{Im}}}


%categories
\newcommand{\Grp}{{\textbf{Grp}}}
\newcommand{\Fld}{{\textbf{Fld}}}
\newcommand{\Sh}{{\textbf{Sh}}}
\newcommand{\Set}{{\textbf{Set}}}
\newcommand{\Mod}{{\textbf{Mod}}}
\newcommand{\PSh}{{\textbf{PSh}}}
\newcommand{\Sch}{{\textbf{Sch}}}
\newcommand{\Top}{{\textbf{Top}}}
\newcommand{\RgSp}{{\textbf{RgSp}}}
\newcommand{\Ring}{{\textbf{Ring}}}


%operatorname -- 3+ letters (lowercase and upper case)
\newcommand{\rad}{{\operatorname{rad}}}
\newcommand{\sep}{{\operatorname{sep}}}
\newcommand{\res}{{\operatorname{res}}}
\newcommand{\sgn}{{\operatorname{sgn}}}
\newcommand{\pre}{{\operatorname{pre}}}
\newcommand{\ord}{{\operatorname{ord}}}
\newcommand{\vol}{{\operatorname{vol}}}
\newcommand{\lcm}{{\operatorname{lcm}}}

\newcommand{\conj}{{\operatorname{conj}}}
\newcommand{\disc}{{\operatorname{disc}}}
\newcommand{\stab}{{\operatorname{stab}}}
\newcommand{\kdim}{{\operatorname{kdim}}}
\newcommand{\diag}{{\operatorname{diag}}}

\newcommand{\trdeg}{\operatorname{trdeg}}
\newcommand{\coker}{{\operatorname{coker}}}
\newcommand{\codim}{{\operatorname{codim}}}

\newcommand{\Ann}{{\operatorname{Ann}}}
\newcommand{\Aut}{{\operatorname{Aut}}}
\newcommand{\Der}{{\operatorname{Der}}}
\newcommand{\Div}{{\operatorname{Div}}}
\newcommand{\Emb}{{\operatorname{Emb}}}
\newcommand{\End}{{\operatorname{End}}}
\newcommand{\Ext}{{\operatorname{Ext}}}
\newcommand{\Gal}{{\operatorname{Gal}}}
\newcommand{\Hom}{{\operatorname{Hom}}}
\newcommand{\CHom}{{\operatorname{CHom}}}
\newcommand{\Ind}{{\operatorname{Ind}}}
\newcommand{\Inn}{{\operatorname{Inn}}}
\newcommand{\Int}{{\operatorname{int}}}
\newcommand{\Irr}{{\operatorname{Irr}}}
\newcommand{\Jac}{{\operatorname{Jac}}}
\newcommand{\Mor}[1]{\operatorname{Mor}(\textbf{#1})}
\newcommand{\Nil}{{\operatorname{Nil}}}
\newcommand{\Obj}{{\operatorname{Obj}}}
\newcommand{\Out}{{\operatorname{Out}}}
\newcommand{\Pic}{{\operatorname{Pic}\,}}
\newcommand{\Rep}{{\operatorname{Rep}}}
\newcommand{\Res}{{\operatorname{Res}}}
\newcommand{\Spl}{{\operatorname{Spl}}}
\newcommand{\Syl}{{\operatorname{Syl}}}
\newcommand{\Sym}{{\operatorname{Sym}}}
\newcommand{\Tor}{{\operatorname{Tor}}}

\newcommand{\Char}{{\operatorname{char}}}
\newcommand{\Frac}{{\operatorname{Frac}}}
\newcommand{\Frob}{{\operatorname{Frob}}}
\newcommand{\Proj}{{\operatorname{Proj}\,}}
\newcommand{\RMod}{{}_R\operatorname{Mod}}
\newcommand{\SExt}{{\mathcal{E}\emph{xt}}}
\newcommand{\SHom}{{\emph{Hom}}}
\newcommand{\Span}{{\operatorname{span}}}
\newcommand{\Spec}{{\operatorname{Spec}\,}}
\newcommand{\Supp}{{\operatorname{Supp}}}
\newcommand{\Tens}[1]{#1\otimes --}
\newcommand{\fHom}[1]{\operatorname{Hom}(#1, --)}

\newcommand{\mSpec}{{\operatorname{mSpec}}}
\newcommand{\cofHom}[1]{\operatorname{Hom}(--, #1)}


% connectors
% \renewcommand{\mid}{\big|}
\newcommand{\sse}{\subseteq}
\newcommand{\subg}{\leqslant}
\newcommand{\supg}{\geqslant}
\newcommand{\ssne}{\subsetneq}
\newcommand{\isosubg}{\lesssim}
\newcommand{\nsse}{\not\subseteq}
\newcommand{\nsubg}{\trianglelefteq}
\newcommand{\nsupg}{\trianglerighteq}
\newcommand{\csubg}{\blacktriangleleft}
\renewcommand{\mid}{\big|}


% Arrows
\newcommand{\rw}{\rightarrow}
\newcommand{\Rw}{\Rightarrow}
\newcommand{\lw}{\leftarrow}
\newcommand{\Lw}{\Leftarrow}
\newcommand{\lrw}{\leftrightarrow}
\newcommand{\LRw}{\Leftrightarrow}
\newcommand{\acton}{\curvearrowright}
\newcommand{\actson}{\curvearrowright}
\newcommand\mapsfrom{\mathrel{\reflectbox{\ensuremath{\mapsto}}}}

% arrows for commutative diagram: create four new angled double-struck arrows
\newcommand\myrot[1]{\mathrel{\rotatebox[origin=c]{#1}{$\Rightarrow$}}}
\newcommand\NEarrow{\myrot{45}}
\newcommand\NWarrow{\myrot{135}}
\newcommand\SWarrow{\myrot{-135}}
\newcommand\SEarrow{\myrot{-45}}


%shorthands
\newcommand{\oln}[1]{{\overline{#1}}}
\renewcommand{\bf}[1]{\textbf{#1}}
\newcommand{\qn}{\quad\newline}


% limits
\newcommand{\Lim}{{\varprojlim}}
\newcommand{\coLim}{{\varinjlim}}
\newcommand{\Dlim}{\lim\limits_{\longrightarrow}}
\newcommand{\coDlim}{\lim\limits_{\longleftarrow}}


%for saturated ideals in the localization section (I don't like these, I think I'll delete)
\newcommand{\LIm}{{}^e }
\newcommand{\LPre}{{}^c }

%  ╔══════════════════════════════════════════════════════════╗
%  ║                    Custom Commands                       ║
%  ╚══════════════════════════════════════════════════════════╝

% swap varphi and phi
\let\temp\phi
\let\phi\varphi
\let\varphi\temp

\newcommand{\placeholder}{\_\_\_\_\_}

\newcommand{\set}[2]{\left\{#1 \ \middle|\ #2\right\}}

%mod should not have the crazy amount of space to its right!
\renewcommand{\mod}{\,\operatorname{mod}\,}

% dcups
\newcommand{\dcup}{\sqcup}
\newcommand{\Dcup}{\coprod}
% \newcommand{\Dcup}{\bigsqcup}

\newcommand{\nbot}{\mathrel{\rlap{$\bot$}\mkern2mu/}}

%a nice large circle
\newcommand{\tikzcircle}[2][red,fill=black]{\tikz[baseline=-0.5ex]\draw[#1,radius=#2] (0,0) circle ;}%

\makeatletter
\newcommand*\bigcdot{\mathpalette\bigcdot@{.5}}
\newcommand*\bigcdot@[2]{\mathbin{\vcenter{\hbox{\scalebox{#2}{$\m@th#1\bullet$}}}}}
\makeatother

%somehow not a default command
\def\rddots{\mathstrut^{.^{.^{.}}}}

%% new delimiter
\DeclarePairedDelimiter{\ceil}{\lceil}{\rceil}


%  ╔══════════════════════════════════════════════════════════╗
%  ║                     environments                         ║
%  ╚══════════════════════════════════════════════════════════╝


%remember the option: breakable

% create tcolor boxes
\newtcbtheorem[number within=section]{thm}{Theorem}%
{
colback=green!5,
colframe=green!35!black,
fonttitle=\bfseries,
label type=theorem
}{th}

\newtcbtheorem[number within=section, use counter from=thm]{lem}{Lemma}%
{
colback=blue!5,
colframe=black!35!black,
fonttitle=\bfseries,
label type=lemma
}{lm}
\newtcbtheorem[number within=section, use counter from=thm]{prop}{Proposition}%
{
colback=white!5,
colframe=white!35!black,
fonttitle=\bfseries,
label type=proposition
}{pr}
\newtcbtheorem[number within=section, use counter from=thm]{cor}{Corollary}%
{colback=blue!5,
colframe=blue!35!black,
fonttitle=\bfseries,
label type=corollary
}{co}
\newtcbtheorem[number within=section, use counter from=thm]{axiom}{Axiom}%
{colback=black!5,
colframe=yellow!35!black,
fonttitle=\bfseries,
label type=axiom
}{ax}
\newtcbtheorem[number within=section, use counter from=thm]{defn}{Definition}%
{colback=red!5,
colframe=red!35!black,
fonttitle=\bfseries,
label type=definition
}{df}


% for <s-cr> to create \item[$\LRw$] instead of \item
\newenvironment{equivEnumerate}
 {\begin{enumerate}
	}{
\end{enumerate}}



%Custom enviroments
 \newenvironment{note}
 {\begin{tcolorbox}[enhanced, sharp corners, colback=white , borderline={0.2pt}{0pt}{black}]
   \begin{quote}\textbf{Note}:
   }{
   \end{quote}\end{tcolorbox}}

 \newenvironment{intuition}
 {\begin{quote}\textbf{Intuition}:
   }{
   \end{quote}}

\newtcbtheorem[number within=section, use counter from=thm]{titledBox}{}{%
  enhanced,
  sharp corners,
  colback=white,
  colframe=black,
  borderline={0.2pt}{0pt}{black},
  left=2mm,
  right=2mm,
  top=2mm,
  bottom=2mm,
  title style={opacity=1},
  attach title to upper,
  before upper={\textbf{\tcbtitletext}\quad},
  coltitle=black,
  fonttitle=\bfseries,
  separator sign={\quad}
}{box}


%better example and proof environments
\def\exampletext{Example} % If English
\colorlet{colexam}{red!55!black} % Global example color


\newtcbtheorem[number within=section, use counter from=thm]{example}{Example}{% \newtcolorbox[use counter=testexample]{testexamplebox}{%
% Example Frame Start
empty,% Empty previously set parameters
title={\exampletext \thetcbcounter #1},% use \thetcbcounter to access the testexample counter text
% Attaching a box requires an overlay
attach boxed title to top left,
   % Ensures proper line breaking in longer titles
   minipage boxed title,
% (boxed title style requires an overlay)
boxed title style={empty,size=minimal,toprule=0pt,top=4pt,left=3mm,overlay={}},
coltitle=colexam,fonttitle=\bfseries,
before=\par\medskip\noindent,parbox=false,boxsep=0pt,left=3mm,right=0mm,top=2pt,breakable,pad at break=0mm,
   before upper=\csname @totalleftmargin\endcsname0pt, % Use instead of parbox=true. This ensures parskip is inherited by box.
% Handles box when it exists on one page only
overlay unbroken={\draw[colexam,line width=.5pt] ([xshift=-0pt]title.north west) -- ([xshift=-0pt]frame.south west); },
% Handles multipage box: first page
overlay first={\draw[colexam,line width=.5pt] ([xshift=-0pt]title.north west) -- ([xshift=-0pt]frame.south west); },
% Handles multipage box: middle page
overlay middle={\draw[colexam,line width=.5pt] ([xshift=-0pt]frame.north west) -- ([xshift=-0pt]frame.south west); },
% Handles multipage box: last page
overlay last={\draw[colexam,line width=.5pt] ([xshift=-0pt]frame.north west) -- ([xshift=-0pt]frame.south west); }%
}{ex}



\def\prooftext{Proof} % If English
\NewDocumentEnvironment{Proof}{ O{} }
{
    \colorlet{colexam}{black!55!black} % Global example color
    \newtcolorbox[number within=section]{testproofbox}{%
        empty,% Empty previously set parameters
        title={\emph{\prooftext} \thetcbcounter: #1},
        attach boxed title to top left,
        minipage boxed title,
        boxed title style={empty,size=minimal,toprule=0pt,top=4pt,left=3mm,overlay={}},
        coltitle=colexam,
        fonttitle=\bfseries,
        before=\par\medskip\noindent,
        parbox=false,
        boxsep=0pt,
        left=3mm,
        right=0mm,
        top=2pt,
        breakable,
        pad at break=0mm,
        before upper=\csname @totalleftmargin\endcsname0pt,
        % Removed height constraints
        % Added flexible width
        width=\dimexpr\textwidth-3mm\relax,
        % Modified overlays
        overlay unbroken={\draw[colexam,line width=.5pt] ([xshift=-0pt]title.north west) -- ([xshift=-0pt]frame.south west); },
        overlay first={\draw[colexam,line width=.5pt] ([xshift=-0pt]title.north west) -- ([xshift=-0pt]frame.south west); },
        overlay middle={\draw[colexam,line width=.5pt] ([xshift=-0pt]frame.north west) -- ([xshift=-0pt]frame.south west); },
        overlay last={\draw[colexam,line width=.5pt] ([xshift=-0pt]frame.north west) -- ([xshift=-0pt]frame.south west); },%
    }
    \begin{testproofbox}
}
{\end{testproofbox}\endlist}

\newcommand{\Coh}{\operatorname{Coh}}

%  ╔══════════════════════════════════════════════════════════╗
%  ║                    for Dynkin diagrams                   ║
%  ╚══════════════════════════════════════════════════════════╝

\def\row#1/#2!{#1_{\IfStrEq{#2}{}{n}{#2}} & \dynkin{#1}{#2}\\}
\newcommand{\tble}[1]{
   \renewcommand*\do[1]{\row##1!}
   \[
      \begin{array}{ll}\docsvlist{#1}\end{array}
   \]
}

%  ╔══════════════════════════════════════════════════════════╗
%  ║                         links                            ║
%  ╚══════════════════════════════════════════════════════════╝

\definecolor{LinkColour}{HTML}{A64A3B} % favourite
% \definecolor{LinkColour}{HTML}{8C3D32} % 2nd place
\hypersetup{
  colorlinks,
  citecolor=black,
  filecolor=magenta,
  linkcolor=LinkColour,
  urlcolor=blue,
}


%  ╔══════════════════════════════════════════════════════════╗
%  ║                          misc                            ║
%  ╚══════════════════════════════════════════════════════════╝

\stackMath %to be able to stack text in math mode
\makeindex

%import options: all, bibliography, book-formatting, booksSetting, customEnvironments, dynkin, hw-formatting, language-setup, newcommands, packages, preamble, proofs, settings, tcolorbox, test.svg,
\begin{document}
\maketitle

\vspace{2em}
This article is a written record of a mathematical journey. It is written for the curious student of algebraic geometry who has stared at the definition of the Fourier-Mukai transform and wondered: \emph{What is this actually doing geometrically?} We often learn the classical Fourier transform in analysis---integrating a function against a complex exponential to extract its frequencies. Yet, when we move to the derived category of coherent sheaves on an abelian variety, the integral turns into a derived pushforward, the exponential turns into a Poincar\'{e} line bundle, and the frequencies become the dual abelian variety.

The goal of this article is to demystify this translation. We will start with a deceptively simple question: \emph{How do we evaluate a function at a point?} From there, we will rebuild the geometry of line bundles, discover the meaning of holonomy, and finally arrive at the derived category with an intuitive understanding of why every piece of the Fourier-Mukai machinery is not just useful, but absolutely necessary. Along the way, we will state and prove the key results with enough rigor that the reader can see the engine behind the intuition.

\qn
\vspace{2em}

%----------------------------------------------------------------------
\section{Notation and Conventions}
%----------------------------------------------------------------------

Throughout this article, $k$ denotes an algebraically closed field (the reader may safely take $k = \C$). By an \emph{abelian variety} over $k$ we mean a complete, connected group variety; when $k = \C$, this is a complex torus $\C^g / \Lambda$ that admits an embedding into projective space.

We write $\Coh(X)$ for the category of coherent sheaves on a smooth projective variety $X$, and $D^b(X) = D^b(\Coh(X))$ for its bounded derived category. For a morphism $f\colon X \to Y$, the symbols $f^*$, $f_*$, $Rf_*$, and $Lf^*$ denote the usual pullback, pushforward, derived pushforward, and derived pullback functors, respectively. We use the convention that $\CF[n]$ denotes the shift of a complex $\CF$ by $n$ places to the left, so that $(\CF[n])^i = \CF^{n+i}$.

If $A$ is an abelian variety of dimension $g$, we denote by $\hat{A} = \Pic^0(A)$ its dual abelian variety, and by $\CP$ the Poincar\'{e} bundle on $A \times \hat{A}$, normalized so that $\CP|_{A \times \{\hat{0}\}} \cong \CO_A$ and $\CP|_{\{0\} \times \hat{A}} \cong \CO_{\hat{A}}$.

%----------------------------------------------------------------------
\section{Redefining ``Evaluation''}
%----------------------------------------------------------------------

Let $A$ be an abelian variety over $k$. Classically, if we want to evaluate a function $f$ at a point $x \in A$, we might think of taking the inner product of $f$ with a Dirac delta distribution:
\[
    f(x) = \int_A f(y)\, \delta_x(y) \, dy.
\]
In algebraic geometry, our ``functions'' are global sections of the structure sheaf, $f \in \Gamma(A, \CO_A)$. The geometric way to think about such a section is as a scheme morphism $f\colon A \to \A_k^1$, and evaluating $f$ at $x$ means composing with the inclusion $\{x\} \hookrightarrow A$ to obtain a point in $\A_k^1$.

However, the Fourier-Mukai transform operates on the derived category $D^b(A)$. The affine line $\A_k^1$ is not a coherent sheaf, and the morphism $A \to \A_k^1$ does not live in $D^b(A)$. We need to reformulate ``evaluation'' entirely in the language of sheaves and their morphisms.

The key shift is from a \emph{geometric} perspective to an \emph{operator} perspective. Instead of viewing $f$ as a map between spaces, we view $f$ as the operator ``multiplication by $f$.'' This is an endomorphism of the structure sheaf:
\[
    f \in \Hom_{\CO_A}(\CO_A, \CO_A).
\]
Our algebraic Dirac delta is the skyscraper sheaf $\CO_x$ at the point $x$. Recall that $\CO_x$ has stalk equal to the residue field $\kappa(x)$ at $x$ and zero elsewhere. There is a canonical surjection $\pi_x\colon \CO_A \twoheadrightarrow \CO_x$ induced by the quotient $\CO_{A,x} \to \CO_{A,x}/\fm_x = \kappa(x)$.

To ``evaluate'' $f$ at $x$ using only operators, we compose:
\[
    \CO_A \xrightarrow{\;\times f\;} \CO_A \xrightarrow{\;\pi_x\;} \CO_x.
\]
Tracing the global section $1 \in \Gamma(A, \CO_A)$ through this composition: the first map sends $1 \mapsto f$, and the second restricts to the residue field, yielding the class of $f$ in $\CO_{A,x}/\fm_x = \kappa(x)$.

\begin{prop}{Evaluation via Hom-Sets}{eval-hom}
There is a canonical isomorphism $\Hom_{\CO_A}(\CO_A, \CO_x) \cong \kappa(x)$, and under this isomorphism the composite $\pi_x \circ (\times f)$ corresponds to the value $f(x) \in \kappa(x)$.
\end{prop}

\begin{Proof}
Any $\CO_A$-module homomorphism $\varphi\colon \CO_A \to \CO_x$ is determined by the image $\varphi(1) \in \Gamma(A, \CO_x) = \kappa(x)$. This gives the isomorphism. For the composite $\pi_x \circ (\times f)$, we have $\varphi(1) = \pi_x(f) = f \bmod \fm_x = f(x)$.
\end{Proof}

This reformulation is perfectly suited for the Fourier-Mukai transform: both $\CO_A$ and $\CO_x$ are objects of $D^b(A)$, and the evaluation data lives in a Hom-space between them. Before we apply the transform, however, we need to understand the geometric arena in which it operates.


%----------------------------------------------------------------------
\section{The Literal Geometry of Line Bundles}
%----------------------------------------------------------------------

To understand the Fourier-Mukai transform, we must understand the space it maps into: the dual abelian variety $\hat{A}$. We are taught that $\hat{A} = \Pic^0(A)$, the moduli space of degree-zero line bundles on $A$. But what does that \emph{look like} geometrically?

Let us temporarily work over $\C$ and consider a compact Riemann surface $X$ of genus $g$. Algebraically, a line bundle is an invertible sheaf $\CL$. By Serre's GAGA principle, this corresponds to a holomorphic line bundle in the classical sense: a complex manifold $L$ equipped with a holomorphic projection $\pi\colon L \to X$ whose fibers are one-dimensional complex vector spaces. Inside $L$ sits a \emph{zero section}, an embedded copy of $X$ selecting the origin in each fiber.

\subsection{Divisors and Winding Numbers}

How does a line bundle acquire its ``twist''? Suppose $s\colon X \to L$ is a holomorphic section that meets the zero section transversally at a single point $P$. This section defines the effective divisor $D = P$, and the associated line bundle is $\CL \cong \CO_X(P)$.

The zero of $s$ at $P$ has a striking topological consequence. Choose a small coordinate disk $U$ centered at $P$, and let $V = X \setminus \{P\}$.

Over $U$, the section $s$ is represented by a holomorphic function vanishing simply at $P$, so in a local coordinate $z$ centered at $P$ it looks like $z \mapsto z$. If we trace the values of $s$ along a small circle $\gamma$ around $P$, the image $s(\gamma)$ winds once around the origin in the fiber. The winding number is $+1$.

Over $V$, the section $s$ is nowhere zero, so the loop $s(\gamma)$ can be contracted to a point in the punctured fiber $\C^*$. From the perspective of $V$, the winding number is $0$.

How can the same loop have winding number $1$ and $0$ simultaneously? It cannot---unless the total space $L$ is itself topologically nontrivial. When we build $L$ by gluing the trivial bundles $U \times \C$ and $V \times \C$ along the overlap $U \cap V$, the transition function $g_{UV}\colon U \cap V \to \C^*$ must account for this discrepancy, contributing a net rotation by $2\pi$. The bundle is \emph{twisted}.

This twist is measured by the first Chern class $c_1(L) \in H^2(X, \Z)$. More precisely, the degree of a line bundle on a compact Riemann surface---the integer obtained by integrating $c_1(L)$ over $X$---equals the number of zeros minus the number of poles of any meromorphic section, counted with multiplicity. This is the content of the classical degree formula.

\begin{prop}{Chern Class and Divisor Degree}{chern-degree}
Let $\CL$ be a line bundle on a compact Riemann surface $X$, and let $s$ be any nonzero meromorphic section of $\CL$ with divisor $\mathrm{div}(s) = \sum n_i P_i$. Then
\[
    \deg(\CL) \;=\; \int_X c_1(\CL) \;=\; \sum_i n_i.
\]
\end{prop}

The zeros and poles of meromorphic sections are therefore a direct readout of the topological twisting built into the total space of the bundle.


%----------------------------------------------------------------------
\section{Flexible vs.\ Rigid: The Birth of $\hat{A}$}
%----------------------------------------------------------------------

We now arrive at a distinction that lies at the heart of the entire theory: the difference between topology and holomorphic structure.

\emph{Topological manifolds are flexible.} They behave like sheets of rubber; you can stretch, bend, and smoothly deform them without changing their essential character. \emph{Complex manifolds are rigid.} The demand that transition functions be holomorphic paints a perfectly conformal grid on every tangent space. You can rotate and scale this grid, but you cannot shear it or stretch it nonuniformly without violating the Cauchy-Riemann equations.

What happens when we examine a line bundle of degree zero ($c_1(\CL) = 0$)?

Topologically, the total winding is zero. In the flexible world of continuous deformations, we can untwist the bundle completely: every degree-zero line bundle on a compact Riemann surface is topologically isomorphic to the trivial cylinder $X \times \C$. By the Chern-Weil theorem, such a bundle admits a \emph{flat connection}---one whose curvature form vanishes identically.

\begin{defn}{Flat Connection and Holonomy}{flat-holonomy}
A \emph{flat connection} on a line bundle $L \to X$ is a connection $\nabla$ whose curvature $F_\nabla = \nabla^2$ vanishes. Given a flat connection and a loop $\gamma$ in $X$ based at a point $p$, parallel transport around $\gamma$ defines a linear automorphism of the fiber $L_p$. Since $L_p \cong \C$, this automorphism is multiplication by a nonzero scalar $e^{i\theta} \in U(1)$. This scalar is the \emph{holonomy} of $\nabla$ along $\gamma$.
\end{defn}

A compact Riemann surface of genus $g$ has $2g$ independent loops generating $H_1(X, \Z) \cong \Z^{2g}$. A flat connection on a topologically trivial line bundle is specified up to gauge equivalence by its holonomies around these generators. Each holonomy is a unit complex number, so the data amounts to a point in $(S^1)^{2g}$.

In the topological world, if we disliked a particular holonomy, we could define a continuous function that
gradually ``un-rotates'' the fiber as we traverse the loop, thereby killing the phase shift. But in the
holomorphic world, functions on a compact Riemann surface must be holomorphic, and the maximum principle
forces every global holomorphic function to be constant. We simply do not have the flexibility needed to gauge away the holonomy. The holomorphic structure locks each phase shift permanently into place.

\begin{prop}{Holomorphic Rigidity of Flat Bundles}{rigid-flat}
Let $L_1$ and $L_2$ be two flat holomorphic line bundles on a compact complex manifold $X$ with the same underlying topological bundle (both trivial). If the holonomies of $L_1$ and $L_2$ differ along any loop in $X$, then $L_1$ and $L_2$ are not holomorphically isomorphic.
\end{prop}

\begin{Proof}
A holomorphic isomorphism $\varphi\colon L_1 \xrightarrow{\sim} L_2$ would be a nowhere-vanishing holomorphic section of $L_1^{-1} \otimes L_2$. This tensor product is again a flat line bundle, and its holonomy along any loop $\gamma$ is the ratio of the holonomies of $L_2$ and $L_1$ along $\gamma$. If this ratio is nontrivial for some $\gamma$, then $L_1^{-1} \otimes L_2$ has nontrivial holonomy and hence admits no nonzero global holomorphic sections (we will prove this vanishing in Proposition~\ref{prop:vanishing-flat} below). Therefore no such isomorphism $\varphi$ exists.
\end{Proof}

Thus, two flat line bundles with different holonomies represent genuinely distinct holomorphic manifolds.

\subsection{The Dual Torus}

Let $A = \C^g / \Lambda$ be a $g$-dimensional complex torus. Its first homology is $H_1(A, \Z) \cong \Z^{2g}$. To specify a flat, degree-zero holomorphic line bundle on $A$, we choose a holonomy for each of the $2g$ fundamental loops. Each choice lives on a circle $S^1 \cong U(1)$, so the moduli space of all such bundles is
\[
    \Hom(\Z^{2g}, U(1)) \;\cong\; (S^1)^{2g},
\]
which is itself a $2g$-dimensional real torus, or equivalently a $g$-dimensional complex torus.

\begin{defn}{Dual Abelian Variety}{dual-av}
The \emph{dual abelian variety} of $A$ is
\[
    \hat{A} \;=\; \Pic^0(A) \;=\; \{ L \in \Pic(A) \mid c_1(L) = 0 \}.
\]
When $A = \C^g / \Lambda$, the dual is $\hat{A} \cong \C^g / \hat{\Lambda}$, where $\hat{\Lambda}$ is the dual lattice defined by the Appell-Humbert theorem. Every point $\alpha \in \hat{A}$ encodes the holonomy data---$2g$ phases---of a specific flat line bundle $L_\alpha$ on $A$.
\end{defn}


%----------------------------------------------------------------------
\section{The Vanishing Theorem for Flat Bundles}
%----------------------------------------------------------------------

Before constructing the transform, we need a fundamental cohomological fact about nontrivial flat line bundles. This result will appear repeatedly in every computation that follows.

\begin{prop}{Vanishing of Cohomology for Nontrivial Flat Bundles}{vanishing-flat}
Let $A$ be an abelian variety of dimension $g$, and let $L_\alpha \in \Pic^0(A)$ be a flat line bundle corresponding to $\alpha \in \hat{A}$. If $\alpha \neq \hat{0}$ (i.e., $L_\alpha$ is nontrivial), then
\[
    H^i(A, L_\alpha) = 0 \quad \text{for all } i \geq 0.
\]
\end{prop}

\begin{Proof}
This is a consequence of the group structure. Let $\mu\colon A \times A \to A$ denote the group law. For any $a \in A$, let $t_a\colon A \to A$ be the translation $x \mapsto x + a$. A defining property of $\Pic^0(A)$ is that $t_a^* L_\alpha \cong L_\alpha$ for all $a \in A$ (line bundles in $\Pic^0$ are precisely the translation-invariant ones). In particular, translation induces an action of the group $A(k)$ on $H^i(A, L_\alpha)$.

On the other hand, the Mumford seesaw theorem (see \cite{mumford-av}, Chapter~II) implies that this action must be the identity on cohomology for any line bundle in $\Pic^0$. The group $A(k)$ acts on the finite-dimensional vector space $H^i(A, L_\alpha)$ through scalar multiplication by a character $\chi\colon A(k) \to k^*$. Since $L_\alpha$ is nontrivial, this character is nontrivial. But a nontrivial scalar acting as the identity on a vector space forces that space to be zero.

More concretely, pick $a \in A$ such that $\chi(a) \neq 1$. Then for any cohomology class $\eta \in H^i(A, L_\alpha)$, we have $\eta = t_a^*\eta = \chi(a) \cdot \eta$, which gives $(\chi(a) - 1)\eta = 0$. Since $\chi(a) \neq 1$, we conclude $\eta = 0$.
\end{Proof}

This vanishing theorem is the geometric heart of ``destructive interference'': a nontrivial flat bundle has phase shifts that cancel out perfectly when integrated over the whole variety.


%----------------------------------------------------------------------
\section{The Universal Kernel: The Poincar\'{e} Bundle}
%----------------------------------------------------------------------

We now have two distinct spaces at hand: the base abelian variety $A$ (the space of ``positions'') and its dual $\hat{A}$ (the space of flat ``frequencies''). To execute the Fourier-Mukai transform, we need a mechanism that couples them.

In the classical Fourier transform on $\R^n$:
\[
    \hat{f}(\xi) = \int_{\R^n} f(x)\, e^{-2\pi i \langle x, \xi \rangle} \, dx,
\]
the engine is the kernel $K(x, \xi) = e^{-2\pi i \langle x, \xi \rangle}$, a function on the product space $\R^n \times \widehat{\R^n}$. Fixing a frequency $\xi$ and slicing, $K(-, \xi)$ is the character associated to that frequency.

In algebraic geometry, the role of this kernel is played by the Poincar\'{e} bundle.

\begin{defn}{Poincar\'{e} Bundle}{poincare}
The \emph{Poincar\'{e} bundle} $\CP$ is a line bundle on the product $A \times \hat{A}$ characterized (up to isomorphism) by the following universal property: for every $\alpha \in \hat{A}$,
\[
    \CP|_{A \times \{\alpha\}} \;\cong\; L_\alpha,
\]
and $\CP$ is normalized so that $\CP|_{\{0\} \times \hat{A}} \cong \CO_{\hat{A}}$.
\end{defn}

The normalization condition pins down $\CP$ uniquely. Without it, we could twist by any line bundle pulled back from $\hat{A}$.

\begin{prop}{Fiber Interpretation of $\CP$}{poincare-fiber}
For any point $(x, \alpha) \in A \times \hat{A}$, the fiber of $\CP$ at $(x, \alpha)$ is canonically isomorphic to the fiber of $L_\alpha$ at $x$:
\[
    \CP_{(x, \alpha)} \;=\; (L_\alpha)_x.
\]
\end{prop}

\begin{Proof}
By definition, $\CP|_{A \times \{\alpha\}} \cong L_\alpha$. The fiber of $\CP$ at a point $(x, \alpha)$ is the fiber of the restricted bundle $\CP|_{A \times \{\alpha\}}$ at $x$, which is $(L_\alpha)_x$.
\end{Proof}

Why must the kernel be a line bundle and not, say, a vector bundle of higher rank? The Fourier-Mukai transform is meant to act as a ``change of basis.'' If $\CP$ were a rank-$r$ bundle with $r > 1$, tensoring with it would multiply ranks by $r$, fundamentally distorting the ``size'' of objects. Because $\CP$ has rank one, tensoring with it acts like multiplication by a scalar phase $e^{i\theta}$: it \emph{twists} the geometry without bloating the data.


%----------------------------------------------------------------------
\section{Why the Derived Category?}
%----------------------------------------------------------------------

Let $\pi_1\colon A \times \hat{A} \to A$ and $\pi_2\colon A \times \hat{A} \to \hat{A}$ be the coordinate projections. We wish to define a functor taking sheaves on $A$ to sheaves on $\hat{A}$ by analogy with the integral transform. The recipe has three steps:
\begin{enumerate}
    \item \emph{Pull up:} $\pi_1^* \CF$ (lift the data to the product).
    \item \emph{Tensor:} $\pi_1^*\CF \otimes \CP$ (multiply by the kernel).
    \item \emph{Push down:} $\pi_{2*}(\pi_1^*\CF \otimes \CP)$ (integrate out the base variable).
\end{enumerate}

If we attempt this in the abelian category $\Coh(A)$, the result is disappointing. The pushforward $\pi_{2*}$ computes only the zeroth direct image---the global sections $H^0$ along each fiber. Since $A$ is a $g$-dimensional projective variety, there are potentially $g$ higher cohomology groups $H^1, \ldots, H^g$ along each fiber, and the naive pushforward discards all of them. The resulting functor is ``lossy'' and has no hope of being an equivalence.

\begin{defn}{Derived Category and Derived Pushforward}{derived-cat}
The \emph{bounded derived category} $D^b(X) = D^b(\Coh(X))$ is the category whose objects are bounded complexes of coherent sheaves on $X$, with morphisms given by localizing at quasi-isomorphisms. The \emph{derived pushforward} $Rf_*$ of a morphism $f\colon X \to Y$ sends a complex $\CF^\bullet$ to its total higher direct image. In particular, the cohomology sheaves of $Rf_*\CF$ are
\[
    \CH^i(Rf_*\CF) \;=\; R^i f_*\CF,
\]
the $i$-th higher direct image of $\CF$. The object $Rf_*\CF$ packages all of $R^0 f_*\CF, R^1 f_*\CF, \ldots$ into a single complex.
\end{defn}

By replacing $\pi_{2*}$ with $R\pi_{2*}$, we retain the full cohomological information of the fibers. No geometric data is lost.

\begin{defn}{Fourier-Mukai Transform}{fm-transform}
The \emph{Fourier-Mukai transform} is the exact functor $\Phi = \Phi_\CP\colon D^b(A) \to D^b(\hat{A})$ defined by
\[
    \Phi(\CF) \;=\; R\pi_{2*}(\pi_1^* \CF \otimes^L \CP),
\]
where $\otimes^L$ is the derived tensor product (which agrees with the ordinary tensor product when one factor is locally free, as $\CP$ is).
\end{defn}

The central result of the theory is:

\begin{thm}{Mukai's Equivalence Theorem}{mukaiEquiv}
The Fourier-Mukai transform $\Phi_\CP\colon D^b(A) \to D^b(\hat{A})$ is an equivalence of triangulated categories. Its quasi-inverse is
\[
    \Psi(\CG) \;=\; R\pi_{1*}(\pi_2^* \CG \otimes^L \CP^\vee) [g],
\]
where $\CP^\vee$ is the dual of the Poincar\'{e} bundle, and $[g]$ denotes the shift by $g$ in the derived category.
\end{thm}

The proof, due to Mukai \cite{mukai}, proceeds by computing $\Psi \circ \Phi$ using the projection formula and base change, reducing the problem to a cohomological computation on $A \times \hat{A} \times A$ involving the ``convolution'' of the Poincar\'{e} bundle with its dual. The key input is the vanishing theorem (Proposition~\ref{prop:vanishing-flat}), which ensures that the convolution collapses to a shifted structure sheaf of the diagonal. We will see echoes of this argument in the explicit computations below.


%----------------------------------------------------------------------
\section{The Grand Finale: Recovering Evaluation}
%----------------------------------------------------------------------

We now return to the question that opened this article. How does the classical evaluation of a function $f(x)$ look through the lens of the Fourier-Mukai transform? We established in Section~2 that evaluation is encoded by the Hom-space $\Hom_A(\CO_A, \CO_x)$. Let us compute the images of both objects under $\Phi$.

\subsection{Transforming the Skyscraper Sheaf $\CO_x$}

\begin{prop}{Transform of a Skyscraper Sheaf}{transform-sky}
Let $x \in A$. Then
\[
    \Phi(\CO_x) \;\cong\; P_x,
\]
where $P_x = \CP|_{\{x\} \times \hat{A}}$ is a line bundle on $\hat{A}$ sitting in degree $0$.
\end{prop}

\begin{Proof}
Pull back: the sheaf $\pi_1^* \CO_x$ is the structure sheaf of the subvariety $\{x\} \times \hat{A} \hookrightarrow A \times \hat{A}$, which we may identify with $\CO_{\{x\} \times \hat{A}}$.

Tensor: $\pi_1^*\CO_x \otimes \CP \cong \CP|_{\{x\} \times \hat{A}} = P_x$ (viewed as a sheaf on $A \times \hat{A}$ supported on $\{x\} \times \hat{A}$).

Push down: we must apply $R\pi_{2*}$ to $P_x$ supported on $\{x\} \times \hat{A}$. The restriction of $\pi_2$ to $\{x\} \times \hat{A}$ is an isomorphism onto $\hat{A}$. In particular, $\pi_2|_{\{x\} \times \hat{A}}$ is a finite morphism of degree one, so all higher direct images vanish:
\[
    R^i \pi_{2*}(P_x) = 0 \quad \text{for } i > 0.
\]
The zeroth direct image is $\pi_{2*}(P_x) \cong P_x$, viewed as a line bundle on $\hat{A}$ via the identification $\{x\} \times \hat{A} \cong \hat{A}$. Therefore $\Phi(\CO_x) = P_x$ concentrated in degree $0$.
\end{Proof}

The point $x$ has been transformed into a global ``character'' line bundle $P_x$ on the dual variety. This is the geometric counterpart of the classical fact that the Fourier transform of a Dirac delta at $x$ is the character $e^{-ix\xi}$.


\subsection{Transforming the Structure Sheaf $\CO_A$}

\begin{prop}{Transform of the Structure Sheaf}{transform-oa}
We have
\[
    \Phi(\CO_A) \;\cong\; \CO_{\hat{0}}[-g],
\]
where $\CO_{\hat{0}}$ is the skyscraper sheaf at the identity of $\hat{A}$, placed in cohomological degree $g$.
\end{prop}

\begin{Proof}
Pull back and tensor: since $\pi_1^*\CO_A = \CO_{A \times \hat{A}}$, we get $\pi_1^*\CO_A \otimes \CP = \CP$.

Push down: we need $R\pi_{2*}\CP$. For each $\alpha \in \hat{A}$, the fiber of $\pi_2$ over $\alpha$ is $A \times \{\alpha\}$, and the restriction of $\CP$ to this fiber is $L_\alpha$. By the theorem on cohomology and base change (see \cite{mumford-av}, \S5), the stalk of $R^i\pi_{2*}\CP$ at $\alpha$ is related to $H^i(A, L_\alpha)$.

If $\alpha \neq \hat{0}$, then $L_\alpha$ is a nontrivial flat bundle, so $H^i(A, L_\alpha) = 0$ for all $i$ by Proposition~\ref{prop:vanishing-flat}. Therefore $R^i\pi_{2*}\CP$ is supported set-theoretically at $\hat{0}$.

If $\alpha = \hat{0}$, then $L_{\hat{0}} = \CO_A$, and we must compute $H^i(A, \CO_A)$. For a $g$-dimensional abelian variety, we have
\[
    H^i(A, \CO_A) \;\cong\; \bigwedge^i H^1(A, \CO_A),
\]
a vector space of dimension $\binom{g}{i}$. In particular, $H^g(A, \CO_A) \cong k$.

Now, using the base change theorem more carefully: since the support of the higher direct images is $\{\hat{0}\}$ and the cohomology is compatible with base change at that point, one can show that $R\pi_{2*}\CP$ is quasi-isomorphic to $\CO_{\hat{0}}$ placed in degree $g$. The essential point is that the Euler characteristic $\chi(L_\alpha)$ is zero for $\alpha \in \Pic^0(A)$ (this is a consequence of the fact that the Hilbert polynomial of a degree-zero bundle on an abelian variety vanishes). When the only nonzero stalk occurs at a single closed point and lives entirely in degree $g$, the complex reduces to $\CO_{\hat{0}}[-g]$.
\end{Proof}

The ``constant function'' $\CO_A$---data spread uniformly across all of $A$---has been compressed into a single shifted point. This is the geometric version of the classical fact that the Fourier transform of the constant function $1$ is a delta function at the origin.


\subsection{Recovering Evaluation via Parseval's Identity}

Because $\Phi$ is an equivalence of triangulated categories, it preserves Hom-spaces. This is the categorical analogue of Parseval's identity in classical Fourier analysis: the ``inner product'' is preserved.

\begin{thm}{Categorical Parseval Identity}{categoricalParsevalIdentity}
For any $\CF, \CG \in D^b(A)$, the Fourier-Mukai transform induces a natural isomorphism
\[
    \Hom_{D^b(A)}(\CF, \CG) \;\cong\; \Hom_{D^b(\hat{A})}(\Phi(\CF), \Phi(\CG)).
\]
\end{thm}

Applying this with $\CF = \CO_A$ and $\CG = \CO_x$:
\begin{align*}
    \Hom_A(\CO_A, \CO_x)
    &\;\cong\; \Hom_{D^b(\hat{A})}(\Phi(\CO_A),\; \Phi(\CO_x)) \\
    &\;\cong\; \Hom_{D^b(\hat{A})}(\CO_{\hat{0}}[-g],\; P_x).
\end{align*}

By the definition of the shift functor, $\Hom_{D^b}(\CF[-g], \CG) = \Hom_{D^b}(\CF, \CG[g])$, so this equals $\Ext^g_{\hat{A}}(\CO_{\hat{0}}, P_x)$.

\begin{lem}{Ext from a Point to a Locally Free Sheaf}{ext-point-lf}
Let $Y$ be a smooth projective variety of dimension $g$, let $y \in Y$ be a closed point, and let $\CE$ be a locally free sheaf on $Y$. Then
\[
    \Ext^g_Y(\CO_y, \CE) \;\cong\; \CE_y \otimes \bigwedge^g T_{Y,y}^*,
\]
where $\CE_y$ is the fiber at $y$ and $T_{Y,y}^*$ is the cotangent space. In particular, when $Y$ is an abelian variety (so $\omega_Y \cong \CO_Y$), we have $\Ext^g_Y(\CO_y, \CE) \cong \CE_y$.
\end{lem}

\begin{Proof}
We compute $\Ext^g$ using the standard Koszul resolution of $\CO_y$. If $y$ has local coordinates $z_1, \ldots, z_g$, the Koszul complex
\[
    0 \to \bigwedge^g \CO_Y^g \to \cdots \to \bigwedge^1 \CO_Y^g \to \CO_Y \to \CO_y \to 0
\]
is a locally free resolution of $\CO_y$. Applying $\CHom(-, \CE)$ and taking the $g$-th cohomology yields
\[
    \Ext^g(\CO_y, \CE) \;\cong\; \CE_y \otimes \det(\fm_y / \fm_y^2)^* \;=\; \CE_y \otimes \bigwedge^g T_{Y,y}^*.
\]
For an abelian variety, the canonical bundle is trivial ($\omega_A \cong \CO_A$), so $\bigwedge^g T^*_{Y,y}$ is a one-dimensional vector space that is canonically trivialized. Hence $\Ext^g(\CO_y, \CE) \cong \CE_y$.
\end{Proof}

Applying Lemma~\ref{lem:ext-point-lf} with $Y = \hat{A}$, $y = \hat{0}$, and $\CE = P_x$:
\[
    \Ext^g_{\hat{A}}(\CO_{\hat{0}}, P_x) \;\cong\; (P_x)_{\hat{0}}.
\]
By the normalization of the Poincar\'{e} bundle, $\CP|_{\{x\} \times \{\hat{0}\}} = (L_{\hat{0}})_x = (\CO_A)_x \cong \kappa(x)$. Therefore $(P_x)_{\hat{0}} \cong \kappa(x)$, and the entire chain of isomorphisms recovers
\[
    \Hom_A(\CO_A, \CO_x) \;\cong\; \kappa(x),
\]
exactly reproducing the evaluation $f(x)$ from Proposition~\ref{prop:eval-hom}.

We have traced a complete path: from the classical evaluation of a function at a point, through the geometry of flat line bundles and holonomy, across the universal Poincar\'{e} kernel, through the homological algebra of derived pushforwards, and back to evaluation on the dual side. The abstract machinery of the

%----------------------------------------------------------------------
\section{Plane Waves to Delta Spikes: Transforming Flat Bundles}
%----------------------------------------------------------------------

We have seen that a point $\CO_x$ transforms into a line bundle $P_x$. What happens when we feed the transform a flat line bundle---a pure ``plane wave''?

In classical Fourier analysis, a pure exponential $e^{2\pi i \alpha x}$ is a single frequency that fills all of space. Its Fourier transform is a Dirac delta at the frequency $\alpha$. We will now watch this phenomenon play out in perfect algebraic detail.

\begin{thm}{Transform of a Flat Line Bundle}{transfromOfFlatLineBundle}
Let $L_\alpha \in \Pic^0(A)$ be the flat line bundle corresponding to $\alpha \in \hat{A}$. Then
\[
    \Phi(L_\alpha) \;\cong\; \CO_{-\alpha}[-g],
\]
i.e., the result is a skyscraper sheaf at the point $-\alpha$, shifted to degree $g$.
\end{thm}

\begin{Proof}
Pull back and tensor: $\pi_1^* L_\alpha \otimes \CP$. Restricting to the slice $A \times \{\beta\}$, we get $L_\alpha \otimes L_\beta \cong L_{\alpha + \beta}$, using the group structure on $\hat{A}$.

Push down: we must compute $H^i(A, L_{\alpha + \beta})$ for each $\beta \in \hat{A}$.

If $\alpha + \beta \neq \hat{0}$ (equivalently, $\beta \neq -\alpha$), the bundle $L_{\alpha+\beta}$ is nontrivial. By Proposition~\ref{prop:vanishing-flat}, all cohomology vanishes. This is the geometric realization of \emph{destructive interference}: the nontrivial phase shifts cancel perfectly upon integration.

If $\beta = -\alpha$, then $L_{\alpha + \beta} = L_{\hat{0}} = \CO_A$. The cohomology of $\CO_A$ concentrates in degree $g$: $H^g(A, \CO_A) \cong k$ (one-dimensional), with all other $H^i(A, \CO_A)$ contributing to lower degrees. However, exactly as in the proof of Proposition~\ref{prop:transform-oa}, the derived pushforward collapses to a skyscraper at $-\alpha$ in degree $g$.

Therefore $\Phi(L_\alpha) \cong \CO_{-\alpha}[-g]$.
\end{Proof}

The global ``plane wave'' $L_\alpha$ has been compressed into an infinitely sharp spike at $-\alpha$. The minus sign is precisely the standard inversion that appears in classical Fourier analysis; it arises from the convention that the Poincar\'{e} bundle encodes $L_\alpha$ on the slice $A \times \{\alpha\}$, so tensoring $L_\alpha$ with $L_\beta$ yields the trivial bundle at $\beta = -\alpha$.


%----------------------------------------------------------------------
\section{Biduality and the Double Transform}
%----------------------------------------------------------------------

We are now in a position to exhibit the most beautiful symmetry of the theory.

\begin{thm}{Biduality}{biduality}
There is a canonical isomorphism $\hat{\hat{A}} \cong A$, identifying the double dual with the original variety.
\end{thm}

The Fourier-Mukai transform provides a constructive proof. Starting with a skyscraper $\CO_x$ on $A$, we apply $\Phi$ twice.

\begin{cor}{Double Transform}{double-transform}
For any $x \in A$,
\[
    \Phi_{\hat{A} \to \hat{\hat{A}}}\bigl(\Phi_{A \to \hat{A}}(\CO_x)\bigr) \;\cong\; \CO_{-x}[-g].
\]
In other words, applying the Fourier-Mukai transform twice returns the original object, mirrored through the group inverse and shifted by $g$.
\end{cor}

\begin{Proof}
By Proposition~\ref{prop:transform-sky}, the first transform gives $\Phi(\CO_x) = P_x$, a line bundle in $\Pic^0(\hat{A})$. The line bundle $P_x$, viewed as a flat bundle on $\hat{A}$, corresponds to a point in $\hat{\hat{A}}$; under the identification $\hat{\hat{A}} \cong A$, this point is $x$ itself (this is the content of the biduality isomorphism at the level of points).

Now apply the second transform (from $\hat{A}$ to $\hat{\hat{A}} \cong A$) to $P_x$. By the preceding theorem on flat bundles, a flat line bundle corresponding to a point $x \in \hat{\hat{A}}$ transforms into a shifted skyscraper at $-x$:
\[
    \Phi(P_x) \;\cong\; \CO_{-x}[-g].
\]
\end{Proof}

In classical Fourier analysis, the most celebrated inversion identity reads $\CF(\CF(f))(x) = f(-x)$: applying the transform twice yields the original function reflected through the origin. Corollary~\ref{cor:double-transform} is the precise algebraic-geometric incarnation of this identity, with the cohomological shift $[-g]$ accounting for the ``cost'' of performing a derived integral over a $g$-dimensional variety.


%----------------------------------------------------------------------
\section{Gaussians and Vector Bundles: Beyond the Extremes}
%----------------------------------------------------------------------

We have explored the two extremes of the transform:
\begin{itemize}
    \item Points ($\CO_x$) transform into flat line bundles ($P_x$).
    \item Flat line bundles ($L_\alpha$) transform into shifted points ($\CO_{-\alpha}[-g]$).
\end{itemize}

What lies in between? In classical analysis, the archetype of a function that is neither a delta spike nor a pure sine wave is a \emph{Gaussian wave packet} $e^{-x^2}$---localized in position but also localized in frequency. One of the signature properties of a Gaussian is that its Fourier transform is again a Gaussian.

The geometric counterpart of a Gaussian is an \emph{ample line bundle}: a bundle $\CL$ with strictly positive first Chern class. Unlike a flat bundle (which has no curvature) or a skyscraper (which is concentrated at a point), an ample bundle carries a nondegenerate ``bump'' of curvature spread over the entire variety.

\begin{thm}{Transform of an Ample Line Bundle (Mukai)}{transfromOnAmpleLineBundles}
Let $L$ be an ample line bundle on $A$ of degree $d = \chi(L) > 0$. Then
\[
    \Phi(L) \;\cong\; \hat{V}[-g],
\]
where $\hat{V}$ is a locally free sheaf (vector bundle) on $\hat{A}$ of rank $d$, placed in cohomological degree $g$.
\end{thm}

This is a consequence of the \emph{index theorem for abelian varieties}. For an ample line bundle $L$ on $A$, the cohomology of $L$ is concentrated in a single degree:

\begin{lem}{Kodaira Vanishing on Abelian Varieties}{kodaira-av}
If $L$ is ample on $A$, then $H^i(A, L) = 0$ for $i \neq 0$, and $\dim H^0(A, L) = \chi(L) = d$.
\end{lem}

When we tensor $L$ with the fibers of $\CP$ and push down, the curvature of $L$ prevents the destructive interference that killed the cohomology in the flat case. The ``bump'' survives integration, producing nonzero output spread across all of $\hat{A}$. The degree $d$---measuring the total curvature---becomes the rank of the resulting vector bundle.

For example, on an elliptic curve ($g = 1$), transforming a line bundle of degree $3$ yields a rank-$3$ vector bundle on the dual curve. Higher curvature input produces higher-rank output: the ``narrower the Gaussian,'' the ``wider its transform.''


%----------------------------------------------------------------------
\section{WIT Sheaves and the Index of a Sheaf}
%----------------------------------------------------------------------

The examples above share a remarkable feature: the Fourier-Mukai transform of each object was concentrated in a single cohomological degree. This is not always the case for an arbitrary object in $D^b(A)$, but it happens often enough to deserve a name.

\begin{defn}{WIT Sheaves}{wit}
A coherent sheaf $\CF$ on $A$ is said to satisfy the \emph{weak index theorem with index $i$} (or is \emph{$\mathrm{WIT}_i$}) if
\[
    R^j\pi_{2*}(\pi_1^*\CF \otimes \CP) = 0 \quad \text{for all } j \neq i.
\]
In this case, $\Phi(\CF) \cong \hat{\CF}[-i]$ where $\hat{\CF} = R^i\pi_{2*}(\pi_1^*\CF \otimes \CP)$ is a single coherent sheaf on $\hat{A}$.
\end{defn}

All the objects we have studied are WIT:
\begin{itemize}
    \item The skyscraper $\CO_x$ is $\mathrm{WIT}_0$, with $\hat{\CO}_x = P_x$.
    \item The structure sheaf $\CO_A$ is $\mathrm{WIT}_g$, with $\hat{\CO}_A = \CO_{\hat{0}}$.
    \item A flat line bundle $L_\alpha$ is $\mathrm{WIT}_g$, with $\hat{L}_\alpha = \CO_{-\alpha}$.
    \item An ample line bundle $L$ of degree $d$ is $\mathrm{WIT}_0$ when $d > 0$ (by Kodaira vanishing), with $\hat{L}$ a vector bundle of rank $d$.
\end{itemize}

Not every sheaf is WIT. For a general coherent sheaf, the transform $\Phi(\CF)$ will be a genuine complex with nonzero cohomology in multiple degrees, and working with it requires the full power of the derived category.


%----------------------------------------------------------------------
\section{The Explicit Dictionary}
%----------------------------------------------------------------------

For readers who want to see exactly how the categorical formula $\Phi(\CF) = R\pi_{2*}(\pi_1^*\CF \otimes \CP)$ aligns with the integral $\int f(x) e^{-ix\xi}\, dx$, the following table provides a precise translation.

\begin{table}[htp]
\centering
\renewcommand{\arraystretch}{1.5}
\begin{tabular}{|p{4.2cm}|p{4.5cm}|p{6cm}|}
\hline
\emph{Classical Fourier Analysis} & \emph{Categorical Fourier-Mukai} & \emph{Geometric Meaning} \\
\hline
Function $f(x)$ & Sheaf $\CF \in D^b(A)$ & The object encoding our data. \\
\hline
Position $x \in \T^n$ & Point $x \in A$ & A coordinate on the base space. \\
\hline
Frequency $\xi \in \Z^n$ & Line bundle $\alpha \in \hat{A}$ & A specific flat twist (holonomy) to be measured. \\
\hline
Kernel $e^{-ix\xi}$ & Poincar\'{e} bundle $\CP$ & The universal coupling between position and frequency. \\
\hline
Multiplication $f(x) \cdot e^{-ix\xi}$ & Tensor $\pi_1^*\CF \otimes \CP$ & Imprinting the phase shift onto the data. \\
\hline
Integration $\int_{\T^n}(\cdots)\,dx$ & Derived pushforward $R\pi_{2*}$ & Summing the data along the fibers. \\
\hline
\end{tabular}
\caption{The Categorical Dictionary.}
\label{tab:dictionary}
\end{table}

This correspondence is not merely a suggestive analogy. Under the Grothendieck-Riemann-Roch theorem, the pushforward functor $\pi_{2*}$ applied at the level of Chern characters becomes literally the integration of differential forms over the fiber. When you apply the Chern character $\ch$ to the Fourier-Mukai transform and pass to cohomology, the result projects down to the classical Fourier integral on the cohomology ring of $A$.


%----------------------------------------------------------------------
\section{What about $\C$? The Limits of Topology}
%----------------------------------------------------------------------

Internalizing this dictionary leads to a natural question: if the Fourier-Mukai transform generalizes the Fourier transform for compact tori, can we use it for the standard Fourier transform on $\C$ (or $\R$)?

The answer is \emph{no}, and the reason is illuminating.

To build the dual variety $\hat{A}$, we relied entirely on \emph{holonomy}: the ability of flat connections to pick up nontrivial phase shifts around the loops of $A$. The complex plane $\C$ is simply connected; every loop contracts to a point. Because $\C$ has no nontrivial loops, every flat line bundle on $\C$ is holomorphically trivial. The putative ``dual space'' $\Pic^0(\C)$ consists of a single point (the trivial bundle), and a Fourier transform into a one-point space carries no information.

\subsection{The Solution: $\CD$-Modules and the Fourier-Laplace Transform}

If topology (loops, holonomy) cannot create a frequency space, we must turn to \emph{calculus}. The correct framework for the Fourier transform on $\C$ (or more generally, on non-compact spaces) is the theory of $\CD$-modules: sheaves of differential equations.

Consider the Weyl algebra $A_1 = k\langle x, \partial_x \rangle / (\partial_x x - x \partial_x - 1)$, the ring of polynomial differential operators in one variable. It is generated by two operations: multiplication by $x$, and differentiation $\partial_x$.

Recall the classical symmetry of the Fourier transform:
\begin{align*}
    \CF(\partial_x f)(\xi) &= i\xi \cdot \CF(f)(\xi), \\
    \CF(x \cdot f)(\xi) &= i\, \partial_\xi \CF(f)(\xi).
\end{align*}
Differentiation becomes multiplication, and multiplication becomes differentiation. In the algebraic theory of $\CD$-modules, we define the \emph{Fourier-Laplace transform} as the automorphism of the Weyl algebra:
\begin{align*}
    x &\longmapsto -\partial_\xi, \\
    \partial_x &\longmapsto \xi.
\end{align*}
This is an algebra isomorphism: both sides satisfy the canonical commutation relation $[\partial, x] = 1$.

Let us test this on our favorite object: the skyscraper at $x = 0$. In $\CD$-module language, the delta distribution $\delta_0$ is characterized by the equation $x \cdot \delta_0 = 0$. Its corresponding $\CD$-module is $M = \CD / \CD x$ (the quotient by the left ideal generated by $x$).

Applying the Fourier-Laplace transform: the generator $x$ maps to $-\partial_\xi$, so $M$ transforms to $\hat{M} = \CD / \CD \partial_\xi$. What is this module? It is the space of ``functions'' annihilated by $\partial_\xi$, i.e., constant functions. A constant function is a plane wave of frequency zero.

We have recovered the classical fact that the Fourier transform of a delta spike is a constant, now through purely algebraic manipulation of differential operators. The two theories---Fourier-Mukai for compact varieties and Fourier-Laplace for $\CD$-modules on noncompact spaces---address different geometric settings but achieve the same conceptual goal: exchanging position and frequency.

%----------------------------------------------------------------------
\section*{Conclusion}
%----------------------------------------------------------------------

The Fourier-Mukai transform is rarely intuitive upon first encounter. The notation of derived functors and pushforwards can obscure the beautiful physical reality underneath.

By stepping back to visualize the flexible topology of divisors, the rigid crystalline structure of holonomy on flat bundles, and the universal evaluation machine of the Poincar\'{e} bundle, the derived category reveals itself not as an abstract playground but as the \emph{exact, necessary environment} required to perfectly preserve the geometry of waves and frequencies across dual spaces.

Every piece of the machinery has a clear geometric role. The Poincar\'{e} bundle is the kernel that couples position to frequency. The derived pushforward is the integral that sums over the base. The shift $[-g]$ is the cohomological price of integration over a $g$-dimensional space. And the vanishing theorem for flat bundles---that single, powerful cancellation---is the geometric engine of constructive and destructive interference that makes the entire transform work.

\end{document}
% \newpage
% \printbibliography
